% 第30章思考题30.1--30.10独立概念解，原书V2 PDF485。
% 30.6源书Cu/Zn核素记号歧义明确标记pending，禁止虚构补齐。
\section*{第30章\quad 核物理：思考题解析}
\paragraph{思考题 30.1}
原子核由近似具有有限空间分布的质子、中子构成，
核子之间强相互作用有短作用程，核物质密度近似饱和。
因而在\emph{尺度估计}时，可把核子画作
半径约 \(1\,\mathrm{fm}\) 的邻近小球，
但量子核子是遵循Pauli统计、
带自旋、具有交换作用及波函数空间延展的多体体系。
\[
\boxed{\text{“核子小球堆积”只能作为密度与几何的粗略图像。}}
\]
它不能准确描述核壳层、衰变、
核反应截面和短程核子关联。

\paragraph{思考题 30.2}
饱和核密度模型中，
每核子平均占有近似相同体积 \(V_0\)，
\[
\frac{4\pi R^3}3\approx AV_0
\quad\Longrightarrow\quad
\boxed{R=r_0A^{1/3}},
\]
其中 \(r_0\sim1.2\,\mathrm{fm}\) 为经验量。
总核物质质量 \(\approx Am_N\)，所以
\[
\boxed{\rho_{\rm nuc}\approx
\frac{Am_N}{(4\pi/3)r_0^3A}
=\frac{3m_N}{4\pi r_0^3}}
\]
与 \(A\) 近似无关。
极轻核的量子表面效应和晕核、
重核形变会偏离简单球模型。

\paragraph{思考题 30.3}
核力作用程短而呈\emph{饱和性}，
核体积内每核子与数量级固定的近邻相互作用，
不像长程库仑作用会与所有 \(A\) 个核子累加。
因此总吸引结合能的体积项近似
\[
\boxed{B_{\rm volume}\simeq a_vA}.
\]
核表面核子缺少部分邻居，
引入负的表面修正
\(-a_sA^{2/3}\)；其余有库仑、
对称和配对等项。
故“\(B\) 正比 \(A\)”仅是主导项，
并非所有核的精确成键规律。

\paragraph{思考题 30.4}
在约1fm量级的核内，
质子与中子均受强核力吸引，
可胜过单个质子对之间的库仑排斥。
中子虽无正电，却提供额外核力和
恰当的中子—质子对称能，
因而有助于稳定部分核素。
当 \(Z\) 很大时，
库仑自能 \(U_C\propto Z^2/R\)
增长比局部核力的体积能更快，
重核更容易发生\(\alpha\)衰变、
自发裂变及其他不稳定过程。
\[
\boxed{\text{稳定性是短程核力、库仑能、
Pauli排斥和核壳效应共同竞争的结果。}}
\]

\paragraph{思考题 30.5}
均匀带电球总电荷 \(Q=Ze\)、半径 \(R\)
的经典静电自能是
\[
U_C=\frac{3Q^2}{20\pi\epsilon_0R}
=\frac35\frac{Z^2e^2}{4\pi\epsilon_0R}.
\]
用 \(R=r_0A^{1/3}\)，于是
\[
\boxed{U_C\simeq a_c\frac{Z^2}{A^{1/3}},
\quad a_c=\frac{3e^2}{20\pi\epsilon_0r_0}
=\frac35\frac{e^2}{4\pi\epsilon_0r_0}}.
\]
取 \(r_0=1.2\,\mathrm{fm}\)，
\(e^2/(4\pi\epsilon_0)=1.440\,\mathrm{MeV\,fm}\)：
\[
\boxed{a_c\simeq0.72\,\mathrm{MeV}}.
\]
结合能中库仑项为
\(-a_cZ(Z-1)/A^{1/3}\) 的写法，
以 \(Z(Z-1)\) 近似排除每个质子的自作用；
连续带电球 \(Z^2\) 是大 \(Z\) 的简化。

\paragraph{思考题 30.6}
用质量数 \(A\) 与核电荷 \(Z\) 守恒。
\[
\boxed{{}^{238}_{92}\mathrm U
\to{}^{234}_{90}\mathrm{Th}
+{}^{4}_{2}\mathrm{He}},
\]
\[
\boxed{{}^{90}_{38}\mathrm{Sr}
\to{}^{90}_{39}\mathrm Y
+e^-+\bar\nu_e}.
\]
对铜至镍的质量数不变、核电荷减1的形式，
若假设发生 \(\beta^+\) 衰变，形式为
\[
\boxed{{}^{A}_{29}\mathrm{Cu}
\to{}^{A}_{28}\mathrm{Ni}
+e^++\nu_e},
\]
也有电子俘获通道
\({}^{A}_{29}\mathrm{Cu}+e^-\to{}^{A}_{28}\mathrm{Ni}
+\nu_e\)；哪一种能发生还取决于质量差。
\textbf{源书歧义：}扫描上最后一项写成
“\({}^{29}\mathrm{Cu}+?\to{}^{29}\mathrm{Zn}\)”，
且给出的左右核素标号似与实际
\(Z_{\rm Cu}=29,\ Z_{\rm Zn}=30\) 不完全一致。
若本意是通常的\(\beta^-\)变换，
规范写法应为
\[
\boxed{{}^{A}_{29}\mathrm{Cu}
\to{}^{A}_{30}\mathrm{Zn}
+e^-+\bar\nu_e},
\]
而原式中“\(+?\)”在\emph{反应物侧}
不能直接填同一个衰变电子以完成该式。
若将它按另一种弱过程（如正电子俘获）
解释，则需要另给能量与中微子细节。
\emph{在未核对权威原版前，本思考题第四小问标记
unresolved\_source\_notation，而不是编造一个唯一答案。}

\paragraph{思考题 30.7}
从\({}^{235}_{92}\mathrm U\) 变为
\({}^{207}_{82}\mathrm{Pb}\)。
\(\alpha\)衰变每次质量数减4，
故
\(235-4N_\alpha=207\)，
\[
\boxed{N_\alpha=7}.
\]
随后质子数由
\(92-2(7)=78\) 增至82，
\(\beta^-\)衰变每次
\(Z\mapsto Z+1\)，因此
\[
\boxed{N_{\beta^-}=82-78=4}.
\]
\(\gamma\)跃迁不改变 \(A,Z\)，
其次数不能由元素守恒推出。

\paragraph{思考题 30.8}
自由核发射\(\gamma\)光子，
因动量守恒必须反冲，近似反冲能
\[
\boxed{E_R=\frac{E_\gamma^2}{2Mc^2}}.
\]
发射光子中心能量比核内部跃迁能少约 \(E_R\)，
静止自由核吸收同样跃迁反而需多约 \(E_R\)，
两者中心失谐约 \(2E_R\)；
若大于自然线宽，发射和吸收的共振重合很差。
在固体晶格中，部分事件把动量转移给整个晶体，
\emph{不激发声子}，等效反冲能极小，
形成无反冲共振的\emph{穆斯堡尔效应}。
严格地说不是“单核绝对不能吸收\(\gamma\)”，
而是普通自由核共振线中心难以精确重合。

\paragraph{思考题 30.9}
截面不是单纯的硬球几何面积，
而是反应概率除以入射通量的\emph{有效面积}。
低能中子德布罗意波长可远大于核半径，
其波动散射、穿透、束缚态和
共振吸收使有效截面显著放大。
例如孤立窄共振附近的Breit--Wigner形式
\[
\boxed{\sigma(E)\propto\frac{\Gamma_{\rm in}
\Gamma_{\rm out}}{(E-E_R)^2+(\Gamma/2)^2}}
\]
可比 \(\pi R^2\) 大许多，
但仍必须满足部分波幺正性界限。
所以不存在
\(\sigma\le\pi R^2\) 的普适硬球上界。

\paragraph{思考题 30.10}
\(Q>0\) 只说明最终体系静质量能差是放能，
不自动说明初态粒子能克服\emph{库仑势垒}
或能满足角动量、选择定则及反应截面要求。
对于\(Q<0\)的吸热反应，
实验室阈能常需大于 \(|Q|\)，
原因是末态产物还须满足动量守恒而具有
共同质心运动；对带电入射粒子，
即使 \(Q>0\) 也可能有显著势垒。
量子隧穿允许低于经典势垒的少量反应，
故不能把势垒当绝对阶跃阈值。
把束流加速至过高能量还会导致
低效率、非目标反应道、散射和热损失，
\[
\boxed{\text{最大反应截面与最高净能量效率
都不必对应于最高入射能量。}}
\]
